% Generated draft: requires mathematical and source review.
\documentclass[11pt]{article}
\usepackage[margin=25mm]{geometry}
\usepackage{fontspec}
\setmainfont{Latin Modern Roman}
\usepackage{amsmath,amssymb,mathtools,unicode-math}
\setmathfont{Latin Modern Math}
\usepackage{enumitem,needspace,xurl,hyperref}
\setlength{\parindent}{0pt}
\setlength{\parskip}{5pt}
\setlength{\emergencystretch}{3em}
\newcommand{\N}{\mathbb{N}}
\newcommand{\Z}{\mathbb{Z}}
\newcommand{\Q}{\mathbb{Q}}
\newcommand{\R}{\mathbb{R}}
\newcommand{\C}{\mathbb{C}}
\newcommand{\F}{\mathbb{F}}
\newcommand{\A}{\mathbb{A}}
\newcommand{\PP}{\mathbb{P}}
\newcommand{\E}{\mathbb{E}}
\newcommand{\eps}{\varepsilon}
\newcommand{\dd}{\,\mathrm{d}}
\begin{document}
\section*{UnsolvedMath ID: 10600009}
\textbf{Source problem code:} catalogue entry
\subsection{Definitions and mathematical statement}
Let \(VB_n\) be the virtual braid group on \(n\) strands. In the usual birack coloring action, a birack has an invertible switch \(S:X^2\to X^2\) satisfying the set-theoretic Yang--Baxter equation and the required birack nondegeneracy axioms. Classical generators act by \(S\) on adjacent coordinates, inverse crossings by \(S^{-1}\), and virtual crossings by the corresponding virtual switch, ordinarily the coordinate flip. This gives an action
\[
 \rho_n:VB_n\longrightarrow\operatorname{Bij}(X^n)
\]
when the defining mixed relations are satisfied. The exact source birack/virtual-switch convention is not included.

The first question is whether there exists a birack whose virtual-braid action is faithful: \(\ker\rho_n=\{1\}\), so distinct braids induce distinct maps. The excerpt does not say whether a single birack must work simultaneously for all strand numbers or a separate choice is allowed for each fixed \(n\). No finiteness of \(X\) is required.

The second question is whether a particular virtual-braid invariant cited in the original survey is faithful. Both citations in that question were lost in extraction, so the invariant cannot be identified or replaced by a different representation.

The source states that the action defined by linear biquandles is not faithful, and says this ``almost certainly'' means the corresponding linear invariants of virtual knots and links are not faithful. The latter is retained as a qualified expectation, not a proved consequence: faithfulness of a braid action and completeness of a closed-knot/link invariant are distinct properties. No unnamed braid invariant is declared equal to a linear-biquandle invariant.
\subsection{Short English statement}
Does some birack distinguish all virtual braids through its action, and is the separately cited braid invariant faithful? The latter citations are missing. The source's expectation about closed virtual-knot invariants remains explicitly conjectural.
\subsection{Sources}
\begin{itemize}
\item Local UnsolvedMath registry, permanent ID 10600009, source code catalogue entry; no proposer or proposal year supplied; updated 2026-08-06.
\item \emph{Unsolved Problems in Virtual Knot Theory and Combinatorial Knot Theory}, revised survey problem-list item 9: \url{https://arxiv.org/abs/1409.2823}. Registry access date 2026-07-29; source-TeX extraction. Provenance reports 60 top-level items with compound items retained. Present primary-literature status and redistribution terms require review.
\item The original statement has empty citation slots for its second invariant question and its closing comparison; no bibliography keys are supplied.
\item The summary lists ``Representations of virtual braids by automorphisms and virtual knot groups'' (Bardakov et al., 2016, arXiv:1603.xxxx); ``Virtual and universal braid groups, their quotients and representations'' (2021, arXiv:2107.xxxx); ``Representations of flat virtual braids by automorphisms of free group'' (2023, arXiv:2306.xxxx); ``Maps from braids to virtual braids and braid representations'' (2022, arXiv:2210.xxxx); and ``The nontrivial kernel of Manturov-Nikonov map from classical braids to virtual braids'' (2026, arXiv:2603.xxxx). All five identifiers contain placeholders and are not valid verifiable arXiv references as supplied; none is presented as a checked citation.
\end{itemize}
\end{document}
